
\subsection{Experiment 3: Inductive Learning of Low-Dimensional Embedding}

The two previous experiments assumed that we have access to the entire graph and sufficient memory to store it. 
However, these assumptions may not hold in scenarios like run-time inference or large swarms navigating complex environments. 
In such cases, the node label prediction task transforms into an inductive learning problem. 
This section demonstrates how the same graph formalism, as applied in experiments 1 and 2, can be applied to inductive learning with minor modifications.

Following the network architecture described in Figure \ref{fig:single_graph_model}, 
we obtain a 4 dimensional latent embedding as the output of the last linear layer, 
and pass it through a fully connected network that forms the classification head. 
The end to end network is trained with cross entropy objective, 
where the true node classes form the labels that are compared against logits predicted by the classification head.
We apply a Dropout of 0.2 at the penultimate layer during training, and a weight decay of 5e-4. 

In addition to the cross-entropy objective, we present results using reconstruction loss alone and 
a combination of reconstruction and cross-entropy loss. 
\cd{If the combined loss produces better results} We hypothesize that reconstruction loss acts as a regularizer, 
preventing the network from overfitting to the classification objective by enforcing the graph topology in the latent space. 
This class-agnostic structure enhances the network's generalization ability, a crucial aspect in the inductive setting, 
where partially observed graphs may introduce noise during training. The notion of combining classification and reconstruction loss 
has previously been explored in the context of image classification in \cite{robert2018hybridnet}. 


In case of ABMs, each simulation can run for different time steps, and can visit multiple different nodes. If we consider each simulation a graph, the graphs can be of varying sizes. This task was trained considering one simulation as a graph, and a batch size of one. This enables us inductively learn graph embeddings without explicitly processing node neighbourhoods and batches, while keeping memory requirements low.

(in case results aren't as good) Results show that the task performance using inductive learners, without explicit batch processing to ensure frequent nodes, as discussed in the Section~\ref{sec:motivateGAT}, 
yields poor classification results. Another factor in the poor classification is the disbalance between classes, which makes nodes in class 3 (Fast/Failure) have poor results. 

(in case they do work)
Results show that the task performance using inductive learner based on GAT, without explicit batch processing, but using frequent nodes, is still able to have comparable performance to 
transductive learning. This tells us that we can do real-time prediction and learning for multi-agent swarms, based on new graphs and nodes as they show up in simulations. 

We also see that the GAT based classifier has similar performance as the GradientBoostingClassifier on the penultimate layer, which further motivates the use 
of GAT for learning the classes from graphs. Using the current framework, any size of graph can be input to the model to further train it and make predictions. 

Training without frequent node filtering (including all nodes) slightly reduces the performance of the classification task, which is expected. 